\documentclass[11pt]{article}
\usepackage[T1]{fontenc}
\usepackage{lmodern}
\usepackage{amsmath,amsthm,amssymb,mathtools}
\usepackage[margin=1.12in]{geometry}
\usepackage[final]{microtype}
\usepackage[hidelinks]{hyperref}
\usepackage{url}
\usepackage{enumitem}
\setlist{leftmargin=*}
\allowdisplaybreaks[2]
\newtheorem{theorem}{Theorem}[section]
\newtheorem{lemma}[theorem]{Lemma}
\newtheorem{proposition}[theorem]{Proposition}
\newtheorem{corollary}[theorem]{Corollary}
\theoremstyle{definition}
\newtheorem{definition}[theorem]{Definition}
\theoremstyle{remark}
\newtheorem{remark}[theorem]{Remark}
\title{Rational points of reduced height\\ in the middle-third Cantor set}
\author{Mira}
\date{19 November 2026}
\begin{document}
\maketitle
\begin{abstract}
Let \(C\) be the middle-third Cantor set. For every \(\varepsilon>0\), we bound the number of distinct rational points of \(C\) of reduced denominator at most \(T\) by \(O_\varepsilon(T^{\log_3 2+\varepsilon})\). This is the counting bound conjectured by Broderick, Fishman, and Reich. We decompose the denominator of a purely periodic point by descending through divisors of its period. At each large step, additive estimates on rational grids count the high-entropy part of the new denominator, and the entropy missing from the other components pays for naming their primes. The resulting local bounds multiply at the original denominator scale. We also show that the count is at least a constant times \(T^{\log_3 2}\log T\exp((\log\log T)^\rho)\) for any fixed \(0<\rho<(32\sqrt e/15)\log_3 2\) and all sufficiently large \(T\). The same gain holds within each fixed-prefix cylinder of depth up to any fixed fraction of \(\log_3 T\), after scaling by its Cantor mass. The lower bound uses common periods built from primes whose predecessors have only small prime factors, and rules out every fixed power of \(\log T\) as a correction to \(T^{\log_3 2}\).
\end{abstract}
\section{Introduction}

Let \(C=\{\sum_{j\ge1} a_j3^{-j}:a_j\in\{0,2\}\}\) and
\(\kappa=\log 2/\log 3\). For \(T\ge2\), put
\[
 N_C(T)=\#\{a/q\in C:(a,q)=1,\ 1\le q\le T\}.
\]
We count distinct fractions, including \(0\) and \(1\). The original
symbolic denominator of a periodic digit expansion can substantially
exceed the reduced denominator. Counting short expansions alone therefore
does not give the required upper bound.

\begin{theorem}\label{thm:main}
For every \(\varepsilon>0\) one has
\[
                 N_C(T)\ll_\varepsilon T^{\kappa+\varepsilon}
                 \qquad(T\ge2).
\]
\end{theorem}

The theorem proves Conjecture~3.3 of Broderick, Fishman, and
Reich~\cite{BFR2011}. Their formulation uses the shells
\(3^j\le q\le3^{j+1}\): the cumulative and shell assertions follow
from each other by restriction and geometric summation, respectively,
with the positive exponent slack rescaled by \(\kappa\). It is the
middle-third case of \cite[Conjecture~1.1]{CVY2026}, not an assertion
for other missing-digit sets. The conjecture grew out of
Mahler's question about rational approximation inside \(C\)~\cite{Mahler1984}.
Every two distinct fractions up to denominator \(T\) are \(T^{-2}\)-separated;
ordinary packing in \(C\) thus gives \(O(T^{2\kappa})\). Theorem~\ref{thm:main}
identifies the power at the scale of the denominator itself.

Rahm, Solomon, Trauthwein, and Weiss proved
\(N_C(T)\gg T^\kappa\log T\)~\cite[Proposition~3.3]{RSTW2020}.
One can see the origin of the logarithm directly. Set \(L=\lfloor\log_3 T\rfloor\).
For \(1\le i<L\), take the expansions \(0.A\overline B\) with prefix
\(A\) of length \(i\), period \(B\) of length \(L-i\), and different final
trits in the two words. The numerator before reduction is
\(A(3^{L-i}-1)+B\), viewing the words as base-three integers. It is
prime to three, so the reduced denominator has exact \(3\)-adic
valuation \(i\), and is at most \(3^L\). For each \(i\) there are
\(2^{L-1}\) points, since different all-\(\{0,2\}\) addresses represent
distinct points. Thus the power of \(T\) in the theorem is sharp,
but the \(\varepsilon\) cannot simply be deleted.

\subsection*{The local counting problem}

First consider denominators \(q\) prime to \(3\). The action of
multiplication by \(3\) on \(a/q\) has order \(n=\operatorname{ord}_q(3)\),
and every iterate stays in \(C\). We descend through divisors \(h\) of
\(n\), retaining the real difference \(\{3^h x\}-x\), where braces
denote fractional part. Its denominator is \(q/(q,3^h-1)\).
If a descent step increases this denominator by a factor \(u\),
then the child differences at fixed \(u\), with the preceding real
difference fixed, sit on a single translate of the order-\(u\) circle
grid. The product of the factors \(u\) and the remaining terminal
factor is \(q\).

Here is the quantitative balance behind the proof. At \(u\asymp U\)
the difference set of the child grid is enlarged to a small-doubling,
tripling-invariant set of size \(U^\beta\), say. Boolean sumset growth
bounds the number of child values by \(U^{\kappa\beta+o(1)}\), with
arbitrarily small power slack. To count the moduli \(u\) attaining that
size, split each into entire prime-power factors. Iterated addition
and rotation show that the factors with substantial digit entropy have
almost full combined projection. A common detector on least common
multiples lets us count those nearly full projections at a fixed
complement. The number of names for the complementary factors then
costs at most \(U^{(1-\beta)/2}\), up to small losses, because an inactive
prime leaves a deficit of roughly the logarithm of that prime. Hence
\[
                  \kappa\beta+\tfrac12(1-\beta)\le\kappa .
\]
The bound is for the sum over moduli and child values at one actual
parent, not a uniform assertion about a single denominator.

Section~\ref{sec:itinerary} constructs the descent and accounts for the
short steps jointly. Sections~\ref{sec:detector} and~\ref{sec:native}
give the common detector, which works even when the independently chosen
grid fibers at a common lag have different offsets. The Boolean estimate
and its least-common-multiple application are in
Sections~\ref{sec:boolean}--\ref{sec:lcm}. Section~\ref{sec:tariff}
pays for the complementary primes, and Section~\ref{sec:assembly}
assembles the counting argument and restores preperiods.

We write \(G_u=u^{-1}\mathbb Z/\mathbb Z\) and use additive notation on
the circle. If \(A\) is a set and \(m\) a nonnegative integer, \(mA\)
denotes its \(m\)-fold sumset while \(m\cdot A=\{ma:a\in A\}\)
denotes its scalar image. All entropies and unmarked logarithms use
natural logarithms. Implied constants may depend on fixed tolerances,
but not on the height or the choices of affine offsets.

\section{A divisor itinerary and its small observations}\label{sec:itinerary}

Fix $0<\eta<1/4$, put $\xi=\eta^3$, and fix a threshold
$P_0\ge2$, independent of the height $T$.  We count nonendpoint
points $x=a/q\in C$ in lowest terms, with $(q,3)=1$ and
$T/2<q\le T$.  We can discard $q=1$.  Tripling modulo one permutes
the points of each such denominator: it shifts their periodic ternary
expansions, and its $n$th iterate is the identity, where
$n=\operatorname{ord}_q(3)$.  We use the following published consequence
of Bourgain's theorem: for every $\delta>0$, $n\ll_\delta q^\delta$
for represented $q$ coprime to $3$ \cite[\S6.2]{RSTW2020}. Thus
$n=T^{o(1)}$ uniformly here.  No bound of logarithmic order is needed.

For $h\mid n$ let
\[
 b(h)=(q,3^h-1),\qquad d(h)=q/b(h),\qquad
 W_h=\{3^h x\}-x.
\]
The last variable is a \emph{real} number in $(-1,1)$, whose class in
$\mathbb R/\mathbb Z$ has exact denominator $d(h)$ (the denominator
of zero is one). Initially $h=n$, $b(n)=q$ and $W_n=0$.
Write $m_{>P_0}$ for the factor of $m$ that remains on removing all
prime powers based at primes $\le P_0$.

\begin{lemma}[The denominator itinerary]\label{lem:itinerary}
There is a deterministic rule depending only on $q,T,\eta,P_0$ for
choosing a chain $n=h_0>\cdots>h_J=h_*$ of divisors, with
$h_*$ supported on primes $\le P_0$. The increments
$u_i=b(h_{i-1})/b(h_i)$ and terminal factor $b=b(h_*)$ satisfy
$b\prod_i u_i=q$. Each retained step is of one of two types:
\begin{enumerate}
\item A large step $h_{i-1}=ph_i$, where $p>P_0$ is prime and
$u_i>T^\eta$. Put $L=(h_{i-1})_{>P_0}$. Of the prime powers
$r^e\parallel b(h_{i-1})$ with
$(\operatorname{ord}_{r^e}3)_{>P_0}<L^\xi$, those not surviving
in full in $b(h_i)$ have product at most $T^\xi$.
\item A group of consecutive small steps, of combined increment at
most $T^{2\eta}$; its terminal and initial lags are both retained.
\end{enumerate}
There are $J\le 3/\eta+1$ steps.  The number of possible labels
consisting of $n$, the retained lags and types, and dyadic size
bins for all $u_i$ and $b$, is $T^{o(1)}$ with parameters fixed.
\end{lemma}
\begin{proof}
At lag $g$ with $L=g_{>P_0}>1$, choose a candidate $p\mid L$ at random
with probabilities $v_p(L)\log p/\log L$. A full component
$r^e\parallel b(g)$ of order $d=\operatorname{ord}_{r^e}3$ is not
entirely in $b(g/p)$ exactly if $d\nmid g/p$. As $d\mid g$, this
requires $v_p(d)=v_p(g)$. Its probability is at most
$\log(d_{>P_0})/\log L$. By linearity the expected sum of
$\log(r^e)$ over lost components with $d_{>P_0}<L^\xi$ is at most
$\xi\log b(g)\le\xi\log T$, even charging partially lost components
in full. Take a minimizer, breaking ties by the smallest $p$, and
repeat until $L=1$.

Keep each raw increment $>T^\eta$ unchanged. In each intervening
run, group consecutive increments until their product first reaches
$T^\eta$, then start again; keep a nonempty trailing group even if
its product is one. A completed group has product at most $T^{2\eta}$.
All increments and the terminal multiply to $q$, by cancellation.
There are at most $1/\eta$ large increments, at most $1/\eta$
completed groups, and at most $1+1/\eta$ trailing groups.
The retained lags are $\le n=T^{o(1)}$, and there are $O(\log(2T))$
bins per slot. All labels are functions of the denominator.
\end{proof}

\begin{lemma}[Fibres at fixed past]\label{lem:trace-fibres}
For a retained step $g\to h$, set $B=(3^g-1)/(3^h-1)$.
At fixed real $W_g$ and exact $u=b(g)/b(h)$, the torus values of
$W_h$ lie in a single translate of $G_u=(\tfrac1u\mathbb Z)/\mathbb Z$,
and $u\mid B$. Each such value has at most two real lifts.
At the terminal, fixing $W_{h_*}$ and $b$ puts $x$ in one translate
of $G_b$.
\end{lemma}
\begin{proof}
Modulo one, $B W_h=W_g$, and $d(h)/(B,d(h))=d(g)$. Thus
$(B,d(h))=u$; $d(g)$ is known from $W_g$, so $d(h)=ud(g)$ is
common to all the candidates. Multiplication by $B$ on $G_{d(h)}$
has kernel $G_u$. At the terminal $q=b\,d(h_*)$ is fixed and the
kernel of $3^{h_*}-1$ on $G_q$ is $G_b$.
\end{proof}

\begin{lemma}[A joint entropy bound for the small observations]\label{lem:small-traces}
Fix a nonempty label of Lemma~\ref{lem:itinerary}, and let $X$ be
uniform on all its points. Write $W_i=W_{h_i}(X)$.
There is finite side information $E$ such that, in natural logarithms,
\[
 H(E)+\sum_{i\ {\rm small}} H(W_i\mid E,W_0,\ldots,W_{i-1})
 \ \le\ O\big(\eta\log(2/\eta)\big)\log T+O_\eta(1),
\]
with constants independent of $P_0$. No conditional distribution on a
real past is required to be invariant under tripling.
\end{lemma}
\begin{proof}
Write $X=\sum_{t\ge0}2w_t3^{-t-1}$ with $w_t\in\{0,1\}$ indexed
periodically modulo $n$. The ternary trits of
$Z_i=(1+W_i)/2$ are precisely $1+w_{t+h_i}-w_t$. Indeed their
series sums to $Z_i$ without any carry. These trits are unique:
$Z_i\in(0,1)$ has denominator prime to $3$, hence is not a triadic
endpoint.

For a small step let $U_i$ be the public upper bin endpoint, taken
$\le2T^{2\eta}$. Given $W_{i-1}$, the first
$M_i=\lceil 2\log_3 U_i\rceil+3$ trits specify $W_i$,
\emph{without} being given $u_i$. In fact the torus difference of
two candidates is killed by $B$ in Lemma~\ref{lem:trace-fibres}
and lies in $G_{\operatorname{lcm}(d,d')}$, where
$d,d'$ are their child denominators. Its order divides
\[
 (B,\operatorname{lcm}(d,d'))
 =\operatorname{lcm}((B,d),(B,d'))
 =\operatorname{lcm}(u,u')\le U_i^2.
\]
Distinct real $Z_i$ are thus separated by at least $1/(2U_i^2)$,
including when the torus values agree.

Let $M$ be the maximal such $M_i$ (zero if there are no small
steps). Say that $F_i(t)$ holds if the bits in the entire class
$t+h_i\mathbb Z$ modulo $n$ are constant.  The events decrease as
$i$ increases, and $F_0(t)$ always holds. For $0\le t<M$ record in
$E$ the bit $w_t$ and the first index $i$ at which $F_i(t)$ fails,
or a marker if it never fails. The indices include the large steps.
Then $H(E)\le M\log(2(J+1))$.

From $W_{i-1}$ one knows the differences
$w_{s+h_{i-1}}-w_s$ for every $s$. On each nonconstant
$h_{i-1}$-cycle they recover all the bits: a nonzero difference
fixes both endpoints, after which one integrates around the cycle.
At a site $t<M_i$ of a small observation, its trit is known if
$F_i(t)$ holds. It is also known if $F_{i-1}(t+h_i)$ fails, because
then the target bit is recovered and the source bit is prepaid.
The remaining event
$F_{i-1}(t+h_i)\setminus F_i(t)$ is known given $E,W_{i-1}$:
the first event is read from the old difference cycle and the
second from the recorded index. On that event only the target bit
is unknown. The original uniform distribution is shift-invariant,
since its label depends only on $q$. With $\theta_i=\Pr(F_i(0))$,
the remaining event has probability $\theta_{i-1}-\theta_i$:
indeed $F_i(t)=F_i(t+h_i)\subseteq F_{i-1}(t+h_i)$.
Subadditivity of conditional entropy and the trit reconstruction
give the upper bound $M_i(\theta_{i-1}-\theta_i)\log 2$ for this
step. Across all small steps the sum is at most $M\log 2$, since
the omitted drops at large steps are nonnegative. Now
$M\le4\eta\log_3 T+O(1)$ and $J=O(1/\eta)$.
\end{proof}

\paragraph{The arithmetic available after conditioning.}
At a large raw step $g=ph\to h$ put $L=g_{>P_0}$.
Every increment $u$ splits into (i) a product of size $\le T^\xi$
from the low components of Lemma~\ref{lem:itinerary}, (ii) components
based at primes dividing $n$ not already counted, and (iii) whole
powers $r^e\parallel q$, $r\nmid 3n$, with
$d_r=\operatorname{ord}_r3\mid g$ and $(d_r)_{>P_0}\ge L^\xi$.
Indeed the kernel of reduction of units modulo $r^e$ to units
modulo $r$ has $r$-power order, whereas
$\operatorname{ord}_{r^e}3\mid n$ for $r^e\mid q$. For $r\nmid n$
this gives equality of the orders at every depth up to $e$; the
whole component survives or disappears. At the terminal the analogous
description has no low factor and $d_r\mid h_*$ is supported on
primes $\le P_0$. For fixed $n$, any products $\le T$ supported on
primes dividing $n$ have only $T^{o(1)}$ names: for each fixed
$v>0$ their number is at most
$T^v\prod_{r\mid n}(1-r^{-v})^{-1}
 \le T^v\exp(O_v(\log(2n)))$; then take $v$ arbitrarily small.

\begin{lemma}[Names of subproducts]\label{lem:itinerary-names}
Fix the large-step or terminal order restrictions in the preceding
paragraph. Let $\mathcal R$ be the set of allowed primes $r\nmid3n$.
For $1/2<s<1$, products of distinct bases in $\mathcal R$ of size
$\le R$ have $O_s(R^s)$ names. For $a>0$ restrict further to
$d_r<r^a$; their base products have $O_{a,v}(R^v)$ names for
every $v>a$. Finally products $\le R$ of powers $r^e$ with $e>A$
have $O_{A,v}(R^v)$ names whenever $(A+1)v>1/2$, for an integer
$A\ge1$. Constants here may also depend on $P_0,\xi$, but not on
the lags or $T$. These bounds hold for arbitrary subproducts, even
when selected using the actual point sets at a fixed past.
\end{lemma}
\begin{proof}
At exact order $d$ there are $O(d)$ primes, since their product
divides $3^d-1$. They lie among $1+jd$ for distinct integers $j\ge1$.
Thus their total reciprocal $s$-mass is
$O_s(d^{-s}\sum_{j\le O(d)} j^{-s})=O_s(d^{1-2s})$,
by taking the smallest possible indices. Under $d<r^a$ their
reciprocal $v$-mass is $O(d^{1-v/a})$. Either estimate is
$O(d^{-c})$ for some fixed $c>0$.
Writing $d=ke$ with $k$ supported on primes $\le P_0$ and
$e\mid L$, $e\ge L^\xi$, the sum over the permissible orders of
$d^{-c}$ at a large step is at most
\[
 \prod_{\ell\le P_0\ {\rm prime}}(1-\ell^{-c})^{-1}
 \,\tau(L)L^{-\xi c}=O_{P_0,\xi,c}(1).
\]
At the terminal only the finite Euler product is needed. The
divisor bound makes the estimate uniform even as $L$ varies.
Consequently the Euler product of $1+r^{-s}$, or $1+r^{-v}$ for
the restricted low-order bases, is bounded. For large exponents,
$\sum_{e>A}r^{-ev}\le r^{-(A+1)v}/(1-2^{-v})$; the bounded
base-mass estimate applies with any
$1/2<s<\min(1,(A+1)v)$. Multiplying the corresponding local
factors and applying $1_{m\le R}\le(R/m)^v$ (or exponent $s$)
proves the assertions. None of these bounds enumerates new lag
paths for the subproducts.
\end{proof}

In particular, for a bounded multiplicity $e$ one may group
remaining unramified components by $e$ and bin their \emph{radical}
products: knowing the radical in that slot determines its $e$th
power. Bins for such subproducts at a fixed past are for counting
only. Adding them to the original label before the entropy
calculation would require a new proof of rotation invariance.

\section{One detector for all grids at a given height}\label{sec:detector}

Put $e(x)=\exp(2\pi i x)$, and write $G_q=(q^{-1}\mathbb Z)/\mathbb Z$
for the subgroup of order $q$ of $\mathbb R/\mathbb Z$. We construct one
test, depending on a height bound $S$ but not on $q$, to be used on every
$G_q$ with $q\le S$. All logarithms without a subscript are natural.

\begin{proposition}\label{prop:detector}
For all sufficiently large real $S$, put $t=\log S$ and
\[
 k=\left\lceil\frac{t}{\log\log t}\right\rceil,\quad
 L=\lceil t(\log t)^4\rceil,\quad
 d=\lceil\log_3(2k+1)\rceil,\qquad
 F_S(x)=\frac1L\sum_{j=1}^L e(3^{dj}x).
\]
Then $P_S=|F_S|^{2k}$ takes values in $[0,1]$, $k=o(\log S)$, and
the detector
\[
 D_S(q)=\sum_{z\in G_q}P_S(z)
 \quad\text{satisfies}\quad
 1\le D_S(q)\le q,\qquad \sum_{q\le S}D_S(q)\le S^{1+o(1)} .
\]
\end{proposition}
\begin{proof}
Expand $P_S(x)=\sum_m w_m e(mx)$. The $w_m$ form the probability
distribution of the difference of two independent sums of $k$ frequencies
drawn uniformly with replacement from $\{3^{dj}:1\le j\le L\}$.
The digits of each sum in base $3^d>k$ are its multiplicities. Thus a
collision requires the second tuple to be a permutation of the first, giving
$w_0\le k!/L^k\le (k/L)^k\le S^{-3}$ for large $S$.
Every other frequency satisfies
$\log(2+|m|)\ll dL+\log(2k)=(\log S)^{1+o(1)}$.

For any fixed $0<s<1$ and $m\ne0$, the number $b_S(m)$ of positive
divisors of $|m|$ at most $S$ satisfies
\[
 b_S(m)\le S^s \prod_{p\mid m}(1-p^{-s})^{-1}
 \le S^s\exp\bigl(O_s((\log(2+|m|))^{1-s})\bigr).
\]
Indeed there are at most $\log_2(2+|m|)$ distinct prime divisors,
the $i$th smallest is at least $i+1$, and
$-\log(1-p^{-s})\ll_s p^{-s}$. For each fixed $s$ the exponential factor
on our frequency range is $S^{o_s(1)}$, so $b_S(m)\le S^{o(1)}$
uniformly there, by taking $s$ arbitrarily small. No estimate for the
unrestricted divisor function at these frequencies is needed.
Orthogonality gives $D_S(q)=q\sum_{q\mid m}w_m$. The zero frequency
contributes at most $S^2w_0$ to the sum over $q$; each other $m$ is
charged at most $S b_S(m)$ before multiplication by $w_m$. This proves
the sum bound. The point $z=0$ contributes one and $P_S\le1$.
\end{proof}

\section{Signed sums of independent Cantor fibers}\label{sec:native}

We use $C$ also for the image of the middle-third Cantor set on the
circle. Multiplication by $3$ maps $C$ into itself. For $H=3^h$ with
$h\ge1$ write $[j]_H=1+H+\cdots+H^{j-1}$. If $Z$ is a circle set,
a centered fiber at modulus $q$ means any subset of
$(Z\cap (v+G_q))-v$, with $v$ any point of the circle. Thus independently
chosen fibers can always be added inside $G_q$. We write $b\cdot x$
or $b\cdot A$ for scalar multiplication and $bA$ for the $b$-fold
additive sumset of $A$ for positive integral $b$.

\begin{theorem}\label{thm:native}
For every positive integer $m$ and $\sigma>0$ there are $p_0,S_0$
such that if $S\ge S_0$, $p>p_0$ is prime, $H=3^h$, and
$q\le S$ divides $[p]_H$, then for any $m$ independently chosen
centered fibers $A_i$ of $(H-1)\cdot C$ at $q$,
\begin{equation}\label{eq:native-mixed}
 \left|\sum_{i=1}^m(A_i-A_i)\right|\le S^\sigma D_S(q).
\end{equation}
For fibers of $C$ itself, the same conclusion holds for all $q\le S$,
without $p$, $H$, or any divisibility condition.
\end{theorem}

The first assertion requires the same $H$ for all factors in one sum,
but no relation between their offsets or their witnesses. In particular
further restrictions on each fiber are allowed. A sum of single signed
copies of such sets, with either sign on each, is up to translation a
subset of the sum in \eqref{eq:native-mixed}: fix one element in each
nonempty set and subtract it. The proof has three steps.

\subsection*{A common frequency multiplier}
For a probability $\omega$ on a finite circle grid write
$\widehat\omega(a)=\sum_z \omega(z)e(az)$ and
$E(\omega)=(\sum_z\omega(z)^2)^{-1}$. Reflection is written
$\widetilde\omega$ and convolution is taken with counting measure.
We first allow an arbitrary fixed number $r$ of input laws $\pi_i$,
one per independently centered fiber, and set
$\lambda=\mathop{*}_{i=1}^r(\pi_i*\widetilde\pi_i)$.

Every probability on $C$ has a Fourier coefficient at some
$g\in\{1,\ldots,7\}$ of modulus at least $1/14$. Indeed the Fej\'er
kernel $\tfrac18|\sum_{a=0}^7e(a(x-1/2))|^2$ is at most
$1/(8\sin^2(\pi(x-1/2)))\le1/2$ on $C$. Its constant coefficient
is one and the absolute sum of its other coefficients is seven. Integrating
proves the assertion. It also applies to the distribution of $3^u X$
for any $C$-valued $X$ and integer $u\ge0$.

For each atom of $\pi_i$ fix a $C$-witness $X_i$ with
$(H-1)X_i-v_i$ equal to that atom modulo one, and choose the $X_i$
independently with the induced laws. Fix any sample time $u=dj$ from
Proposition~\ref{prop:detector}. For each $0\le n<M$ there are
$g_i(n)\in\{1,\ldots,7\}$ for which
\[
 \left|\prod_i \mathbb E e(g_i(n)3^u H^n X_i)\right|\ge a,
 \qquad a=14^{-r}.
\]
Pigeonhole the tuple and the angular octant of this product: on a set
$I$ of at least $M/(8\cdot7^r)$ times the tuple $(g_i)$ and octant
are fixed. Put $V_n=\prod_i e(g_i3^uH^n X_i)$ on the product probability
space. Independence and projection onto the central ray of the octant
give $|\mathbb E\sum_{n\in I}V_n|\ge a|I|/2$. Hence
\[
 \sum_{n,n'\in I}\operatorname{Re}\mathbb E(V_n\overline{V_{n'}})
 = \mathbb E\left|\sum_{n\in I}V_n\right|^2\ge a^2|I|^2/4.
\]
Take $M=M(r)$ large enough that $M/(8\cdot7^r)>8/a^2$. Removing
the diagonal, whose contribution is $|I|$, yields a pair $n>n'$
for which $|\mathbb E(V_n\overline{V_{n'}})|\ge a^2/8$. In other
words,
\begin{equation}\label{eq:native-mult}
 \prod_i |\widehat\pi_i(g_i3^u b)|\ge a^2/8,
 \qquad b=H^{n'}[n-n']_H .
\end{equation}
Here recentering contributes only constant phases. Crucially the
possibly large $b$ is the \emph{same} for every $i$. If $p>M$,
multiplication by $b$ permutes $G_q$: $H$ is prime to $[p]_H$, and
$\gcd([l]_H,[p]_H)=[\gcd(l,p)]_H=1$ for $1\le l<M$.
The identity follows by dividing
$\gcd(H^l-1,H^p-1)=H^{\gcd(l,p)}-1$ by $H-1$.
For fibers of $C$ itself one simply takes $b=1$ and applies the
seven-harmonic assertion to the genuine $3^u X_i$, before recentering;
no second moment or unit condition is needed in that case.

Pigeonhole over the $L$ sample times. One tuple $(g_i)$ and one pair
$(n,n')$ (no pair in the case of $C$) work at a fraction
$c_r>0$ of the times. Write $\mu_i=(g_i)_*\pi_i$ for pushforward
under scalar multiplication, $\nu=\mathop{*}_i
(\mu_i*\widetilde\mu_i)$ and $\nu_j=\nu^{*2^j}$ for $j\ge0$.
All Fourier coefficients of \(\nu\) and its convolution powers are nonnegative. By
\eqref{eq:native-mult}, or its $b=1$ version, for each fixed $j$
there is a constant $c_{r,j}>0$ such that
\[
 \sum_{z\in G_q} ((b)_*\nu_j)(z)F_S(z)\ge c_{r,j}.
\]
The expression on the left is real by symmetry. Jensen and
Cauchy--Schwarz give
\[
 c_{r,j}^{2k}\le \sum_z ((b)_*\nu_j)(z) P_S(z)
 \le \|\nu_j\|_2 D_S(q)^{1/2}
\]
since $P_S^2\le P_S$ and $b$ acts bijectively. Consequently
\begin{equation}\label{eq:native-scaled}
 E(\nu_j)\le D_S(q)S^{o_{r,j}(1)}.
\end{equation}
This error is uniform in the laws, $h,p,q$, and offsets once $p>M(r)$,
because $k=o(\log S)$.

\subsection*{Removing the bounded individual multipliers}
We use the weighted Balog--Szemer\'edi--Gowers lemma
\cite[Lemma~1]{TaoBSG}. In the counting-measure normalization used here,
if a symmetric probability $\rho$ on a finite abelian group satisfies
$\|\rho*\rho\|_2>K^{-1}\|\rho\|_2$, $K\ge2$, then there is a
symmetric set $B$ containing zero, with $B+B$ covered by $K^{O(1)}$
translates of $B$, $|B|\le K^{O(1)}E(\rho)$, and
$\rho(v+B)\ge K^{-O(1)}$ for some $v$. All such implied powers
are absolute. In particular $|s B|\le K^{O(s)}|B|$ for positive
integers $s$, by repeatedly using the covering.

Fix $\delta>0$, let $K=S^\delta$, and $J=\lceil2/\delta\rceil$.
All probability norms on $G_q$ lie between $S^{-1/2}$ and one.
Therefore for some $0\le j<J$ we have
$\|\nu_{j+1}\|_2>K^{-1}\|\nu_j\|_2$. Apply the lemma with
$\rho=\nu_j$. By \eqref{eq:native-scaled},
$|B|\le K^{O(1)}D_S(q)S^{o_{r,\delta}(1)}$. Each $\mu_i$ is
a convolution factor of $\nu_j$, so averaging over the other factors
gives a translate $v_i'+B$ on which $\mu_i$ has mass $\ge K^{-O(1)}$.
Thus the restriction of $\pi_i$ to
$S_i=\{x:g_i x\in v_i'+B\}$ has at least this mass.

Put $b_0=\operatorname{lcm}(1,\ldots,7)$. On the event where the
two independent samples of each $\pi_i$ belong to $S_i$, their signed
sum, multiplied by $b_0$, belongs to $(2r b_0)B$. In fact the
translate cancels for each difference and $b_0/g_i$ is a positive
integer. Multiplication by $b_0$ has kernel of size at most $b_0$ on
$G_q$. An event of probability $\ge K^{-O(r)}$ for $\lambda$ is
therefore supported on at most $b_0K^{O(r b_0)}|B|$ points. By
Cauchy--Schwarz and $b_0$ absolute, this proves
\begin{equation}\label{eq:native-effective}
 E(\lambda)\le D_S(q) S^{O_r(\delta)+o_{r,\delta}(1)}.
\end{equation}
Neither concentration of the full supports nor a factorization of their
uniform sumset law has been inferred.

\subsection*{From effective size to full sumsets}
Suppose \eqref{eq:native-mixed} fails for a fixed $m,\sigma$ with
no admissible thresholds. There is then a sequence of violations with
$S\to\infty$ and $p\to\infty$ (only $S\to\infty$ in the terminal
case). At each index keep $H,p,q,S$ fixed and let $M_r$ be the
supremum of $E(\lambda)$ over all allowed $r$-tuples of laws on
independent fibers for these same parameters. Pass to a single
subsequence on which $\log_S M_r\to e_r$ for every fixed $r$ and
$\log_S D_S(q)\to\ell$; all these exponents are in $[0,1]$.
Singleton laws exist. By \eqref{eq:native-effective}, first with
$r,\delta$ fixed and then $\delta\downarrow0$, we have
$e_*:=\sup_r e_r\le\ell$. This passage uses only fixed $r$: the
threshold $p>M(r)$ need not hold simultaneously at an individual
index for all $r$.

Choose $\gamma>0$, then one fixed $r$ with $e_r>e_*-\gamma$, and
laws giving $E(\lambda)\ge M_r/2$. Now choose a new $\delta>0$,
allowed to depend on $r$. At a plateau for the unscaled powers
$\lambda^{*2^j}$, $0\le j<J$, apply the same weighted lemma.
Since $E(\lambda^{*2^j})\le M_{r2^j}$ and there are only finitely
many possible $j$, the symmetric set $B'$ it gives has
\[
 |B'|\le S^{e_*+O(\delta)+o(1)},\qquad
 B'+B'\text{ covered by }S^{O(\delta)}\text{ translates of }B'.
\]
As before, for each $i$ restrict $\pi_i$ to an appropriate
translate of $B'$ carrying mass $\ge S^{-O(\delta)}$, normalizing
to obtain $\pi_i'$. The signed convolution $\lambda'$ of the
restricted laws is supported on $G=(2r)B'$. Pointwise domination
$\pi_i'\le S^{O(\delta)}\pi_i$ gives
$\lambda'\le S^{O(r\delta)}\lambda$ and hence
\[
 E(\lambda')\ge S^{e_*-\gamma-O(r\delta)-o(1)}.
\]

For each nonempty centered fiber $A$, take $Y\subset A$ maximal
with pairwise disjoint translates $y+G$. For the uniform law $u_Y$,
disjointness yields $E(\lambda'*u_Y)=|Y|E(\lambda')$ exactly.
Further convolution with $\widetilde u_Y$ cannot decrease effective
size, by the triangle inequality for convex combinations of translated
laws. The resulting law is allowed in $M_{r+1}$. Therefore, uniformly
for every such $A$,
\[
 |Y|\le S^{\gamma+O(r\delta)+o(1)},\qquad A\subset Y+(G-G).
\]
The inclusion is maximality. Apply it separately to the $m$ fibers of
the purported violation, using the same $G$ for them all. The total
number of choices of the $2m$ centers is at most
$S^{2m\gamma+O_m(r\delta)+o(1)}$; the remaining summands belong
to $(8mr)B'$, of size $\le S^{e_*+O_m(r\delta)+o(1)}$.
Thus the limsup of the log-base-$S$ cardinalities is at most
$e_*+2m\gamma+O_m(r\delta)$. First take $\delta\downarrow0$ for
this fixed $r$, then $\gamma\downarrow0$. The result is at most
$e_*\le\ell$, whereas the violations give a liminf at least
$\ell+\sigma$. This proves Theorem~\ref{thm:native}.

\section{Boolean growth at a fixed real displacement}\label{sec:boolean}

Fix a prime drop from lag \(ph\) to lag \(h\), and write \(H=3^h\),
\(M=[p]_H=1+H+\cdots+H^{p-1}\). We work with actual points
\(x\in C\cap(0,1)\) of reduced denominator \(q\) prime to \(3\);
the two integer endpoints can be counted separately. Fix the signed
\emph{real} number \(W=\{H^p x\}-x\), where braces denote fractional
part. Further restrictions on the possible witnesses are allowed.
Write \(d\) for the denominator of \(W\) (with denominator \(1\) for
zero). For each integer \(u\) let \(A_u\) consist of the distinct
circle values \(w=(H-1)x\bmod 1\) of exact denominator \(ud\).
In fact the denominator of each such \(w\) is a multiple of \(d\):
\(Mw=W\bmod 1\), so \(u=\gcd(M,ud)\). In particular \(u\mid M\);
two such values differ in the kernel of multiplication by \(M\) on
\(G_{ud}=(ud)^{-1}\mathbb Z/\mathbb Z\), which is \(G_u\). Thus
\(A_u\) is contained in one affine \(G_u\)-grid. The real child
\(\{Hx\}-x\in(-1,1)\) has at most two values above any given \(w\).

\begin{proposition}[Growth of a conditional winding]\label{prop:bool-winding}
For every nonempty \(A_u\) as above,
\[
 |A_u|^{1/\kappa}
   \ll p^{3/2}u^{1/(p-1)}\,|A_u-A_u|,
 \qquad \kappa=\log 2/\log 3.
\]
The implied constant is absolute, independent of \(h\) and of the
restrictions on the witnesses.
\end{proposition}

\begin{proof}
Select one witness \(x=x(w)\) for each \(w\). Set
\[
 m=1+\Big\lceil \log_3(4\sqrt p)+\frac{\log_3 u}{p-1}\Big\rceil,
 \qquad x_j=\{H^j x\}\quad(0\le j\le p).
\]
The points \(x_j\) have their unique ternary expansions in the
digits \(0,2\): multiplication by \(H\) removes \(h\) digits, and
their reduced denominators remain prime to \(3\) and exceed one.
Form a Boolean vector \(F(w)\) by recording the first \(m\) ternary
digits, divided by \(2\), at each \(x_j\), \(0\le j<p\).

We first prove injectivity. Suppose \(F(w)=F(w')\), and use primes
on \(x_j\) for the second witness. Put \(v_j=x_j-x'_j\); the prefix
bound gives \(|v_j|\le 3^{-m}\) for \(j<p\). The equality
\(x_p-x_0=x'_p-x'_0=W\) gives \(v_p=v_0\) as real numbers.
Consequently \(t=(v_{j+1}-v_j)_{j=0}^{p-1}\) lies in the real
hyperplane \(E=\{y:\sum_j y_j=0\}\) and has Euclidean norm
\(\le 2\sqrt p\,3^{-m}\). Modulo \(\mathbb Z^p\), its coordinates
are \((H^j(w-w'))_{0\le j<p}\).

Consider the lattice
\[
 \Lambda=\{y\in E:y\bmod\mathbb Z^p
             =(H^j z)_{0\le j<p}\text{ for some }z\in G_u\}.
\]
It contains \(E\cap\mathbb Z^p\) with index at most \(u\), by the
first coordinate of the reduction map. The latter integer lattice
has covolume \(\sqrt p\) in \(E\) (use the basis \(e_j-e_{p-1}\),
\(0\le j<p-1\)). Thus \(\operatorname{covol}(\Lambda)\ge 1/u\).
Since \(u\mid M\), \(H^p z=z\) on \(G_u\); cyclic coordinate shift
preserves \(\Lambda\). If \(w\ne w'\), then \(t\ne0\) is a rational
vector in \(\Lambda\). The shift \(R\) on \(E\cap\mathbb Q^p\)
satisfies \(1+R+\cdots+R^{p-1}=0\). As
\(\Phi_p(X)=1+\cdots+X^{p-1}\) is irreducible over \(\mathbb Q\),
the minimal annihilating polynomial of this nonzero \(t\) is
\(\Phi_p\). Hence \(t,Rt,\ldots,R^{p-2}t\) span a full-rank
sublattice, of covolume at least \(1/u\). Hadamard's inequality
instead bounds their volume by \((2\sqrt p\,3^{-m})^{p-1}<1/u\).
This proves injectivity. Fixing the prior only modulo one would
not give the real equality \(v_p=v_0\) used here.

Next fix a single circle difference \(z=w-w'\). Put
\(e_j=(x_j-x'_j)/2\in(-1/2,1/2)\).
The integers removed when passing from \(H^j x\) to \(x_j\)
have only even ternary digits and are even; the same holds for
\(x'\). Thus \(e_j\equiv H^j e_0\pmod1\).
Writing \(a=(H-1)e_0\bmod1\), we have \(2a=z\), with at most two
choices of \(a\), and
\[
 e_j\equiv e_0+[j]_H a\pmod1,\qquad [0]_H=0.
\]
The ternary digits of \(1/2+e_j\) are one plus the successive
differences of the bits of \(x_j,x'_j\), by the series
\(1/2=\sum_{\ell\ge1}3^{-\ell}\). This number belongs to \((0,1)\)
and has denominator prime to \(3\), so it is not a ternary interval
endpoint. For a fixed \(a\), as the single parameter \(e_0\bmod1\)
varies, the vector of length-\(m\) prefixes of all \(1/2+e_j\) is
constant between at most \(p3^m\) endpoints: these are translates
of the \(3^m\) endpoints of the ternary partition, one partition for
each \(j\). No witness lies on them. The total number of differences
\(F(w)-F(w')\) over this \(z\) is therefore at most \(2p3^m\).

For \(E_1,E_2\subset\{0,1\}^s\) the integer-cube sumset inequality
\cite[Corollary~1.11]{BIMP} gives
\(|E_1+E_2|\ge (|E_1||E_2|)^{(\log_2 3)/2}\). Apply it to
\(F(A_u)\) and its translate-reflection \(\mathbf1-F(A_u)\). Then
\[
 |A_u|^{1/\kappa}\le |F(A_u)-F(A_u)|
     \le 2p3^m |A_u-A_u|.
\]
The definition of \(m\) gives the asserted coefficient.
\end{proof}

\paragraph{The terminal fiber.}
At the final lag \(h\), fix \(W=\{3^h x\}-x\) of denominator \(d\).
For a fixed \(u=\gcd(q,3^h-1)\), the relation
\((3^h-1)x=W\bmod1\) gives \(q=ud\). Its solutions in \(G_q\)
lie in one coset of \(G_u\). Now let \(A_u\) be the set of genuine
points \(x\) on this fiber, read modulo one. We have the stronger
bound
\begin{equation}\label{eq:bool-terminal}
 |A_u|^{1/\kappa}\le 2|A_u-A_u|.
\end{equation}
Indeed take \(m\) with \(3^{-m}<1/u\). Distinct points of one affine
\(G_u\)-grid cannot share their first \(m\) digits. For each fixed
circle difference there are at most two real differences \(x-x'\)
in \((-1,1)\); each specifies one Boolean prefix difference by the
ternary expansion of \(1/2+(x-x')/2\), as above. The same cube
inequality proves \eqref{eq:bool-terminal}.

\section{Stabilization and the quotient estimate}\label{sec:lcm}

Work at one fixed parent, either a drop or a terminal fiber as
above, with fixed common period \(n=\operatorname{ord}_q(3)\)
and one bin \(U/2<u\le U\). The reduced denominators \(q\le T\)
of different witnesses need not coincide. In both cases \(u\mid q\),
so \(\operatorname{ord}_u(3)\mid n\) (take the period to be \(1\)
if \(u=1\)); the represented-period estimate gives \(n=T^{o(1)}\)
uniformly. Discard empty \(A_u\). We distinguish \(mS\), the sum
of \(m\) copies of a set, from \(m\cdot S\), its image under scalar
multiplication. Define the symmetric, exactly tripling-invariant set
\[
 S_u=\bigcup_{0\le t<\operatorname{ord}_u(3)}3^t\cdot(A_u-A_u)
       \subset G_u .
\]
It contains zero. No invariance of \(A_u\) is claimed. For every
\(\theta>0\) choose an integer \(a>1/\theta\). For some
\(j\in\{1,10,\ldots,10^{a-1}\}\) we have
\(|10jS_u|\le U^\theta |jS_u|\): otherwise iteration would give
\(|10^a S_u|>U\). Put \(J=10^{a-1}\), and let \(B_u=4jS_u\) for
one such \(j\). The set \(2B_u=8jS_u\) is covered by at most
\(U^\theta\) translates of \(B_u\). Indeed, take a maximal family
of disjoint translates \(y+jS_u,\ y\in 8jS_u\); all lie in
\(9jS_u\subset10jS_u\). Each unchosen translate meets a chosen
one, so each \(y\) is in a chosen center plus \(2jS_u\subset B_u\).
The set \(B_u\) remains symmetric and exactly tripling-invariant,
and contains \(A_u-A_u\). We obtain, at a drop with \(p>P_0\),
\begin{equation}\label{eq:bool-container}
 |A_u|\ll T^{o(1)} U^{\kappa/(P_0-1)} |B_u|^\kappa ,
\end{equation}
using \(p\le n\); at the terminal the factor preceding
\(|B_u|^\kappa\) can be replaced by \(2^\kappa\).

\begin{lemma}[Quotient inventory at fixed complement]\label{lem:lcm}
Fix \(\theta,\sigma'>0\) and \(\eta>0\), and let
\(T^\eta\le U\le 2T\). For each integer \(1\le K\le U\), put
\(V=U/K\). At either kind of parent, for any subcollection of its
allowed indices,
\[
 \sum_{D:\,KD\ {\rm in\ the\ subcollection}} |K\cdot B_{KD}|
     \ll V U^{\sigma'} T^{o(1)} .
\]
For a drop the bound requires \(p>P_0(\theta,\sigma')\);
there is no time-prime condition at the terminal. The error term
is uniform in the parent, \(K\), the subcollection and the chosen
members of the finite menu. No coprimality of \(K,D\) is assumed.
\end{lemma}

\begin{proof}
Choose a fixed integer \(k>2/\sigma'\). For an ordered tuple of
admissible \(D_i\le V\), write
\[
 L=\operatorname{lcm}(D_1,\ldots,D_k),\qquad
 Q=K L=\operatorname{lcm}(KD_1,\ldots,KD_k).
\]
The identity follows prime by prime, including at primes dividing
\(K\). In the drop case \(Q\mid[p]_H\), for the \emph{same}
\(H,p\) throughout this tuple. Each constituent of
\(\sum_i B_{KD_i}\) can be expanded using the defining union of
rotations. For a choice of rotations the result is a sum of at
most \(4Jk\) sets of the form \(X-X\), where \(X=3^t\cdot A_{KD_i}\)
lies on one affine \(G_Q\)-grid. Indeed \(G_{KD_i}\subset G_Q\);
also \(3^t(H-1)x=(H-1)\{3^t x\}\bmod1\) with
\(\{3^t x\}\in C\). At the terminal \(X\) consists directly of
genuine shifted points \(\{3^t x\}\in C\). Thus all the difference
sets are of the kind tested by the common mixed-fiber estimate.
Their grid offsets may vary independently; the copies need not
share a realized prior. Since taking differences removes every
individual recentering, there is no cost for aligning offsets.

There are at most \(n^{4Jk}\) choices of the rotations. Let \(S\)
be a dyadic detector height with \(Q\le S\). For any fixed
\(\delta>0\), Theorem~\ref{thm:native},
at the finitely many orders up to \(4Jk\), gives
\begin{equation}\label{eq:lift-lcm}
 \left|\sum_{i=1}^k K\cdot B_{KD_i}\right|
 \le \left|\sum_{i=1}^k B_{KD_i}\right|
 \ll n^{4Jk} S^\delta D_S(Q).
\end{equation}
Here \(D_S(Q)\) is the common full-grid detector of
Proposition~\ref{prop:detector}, whose sum over
\(Q\le S\) is \(\le S^{1+o(1)}\). The first inequality is the
cardinality bound for the image under the \emph{single} scalar
\(K\); it uses no division by \(H-1\). At the terminal we use the
true-\(C\) assertion of the same theorem.

On the output side \(K\cdot B_{KD_i}\subset G_{D_i}\).
Addition from \(\prod_i G_{D_i}\) onto \(G_L\) is surjective:
its image is generated by the \(1/D_i\), and
\(\gcd_i(L/D_i)=1\). Each fiber has exactly \(\prod_i D_i/L\)
points. Restricting to the indicated subsets and applying
\eqref{eq:lift-lcm} gives
\[
 \prod_i |K\cdot B_{KD_i}|
       \ll \frac{V^k}{L}\,n^{4Jk} S^\delta D_S(KL).
\]
There are at most \(\tau(L)^k=T^{o(1)}\) ordered tuples for each
\(L\), by the divisor bound since \(L\le V^k\le(2T)^k\).
Group by \(W\le L<2W\) for dyadic \(W\ge1\), taking a \emph{single}
dyadic \(S\) with \(2KW\le S<4KW\) for that shell. Nonempty shells
have \(S\ge Q>U/2\); also \(S\le 4 K V^k\le 4U^k\). Since \(KL\)
is different for each \(L\), the complete detector budget, even
though it includes many unused moduli, gives
\[
 \sum_{W\le L<2W}D_S(KL)\le S^{1+o(1)}.
\]
After expansion of the \(k\)-th power of the desired sum, the shell
therefore contributes at most
\[
 V^k T^{o(1)}\,\frac{S^{1+\delta}}{W}
   \ll_{k,\delta} V^k K U^{k\delta} T^{o(1)} .
\]
The subpower factor includes \(n^{4Jk}\) and the uniform detector
budget error. There are \(O_k(\log(2T))\) shells. Taking a \(k\)-th
root costs \(K^{1/k}\le U^{1/k}\), not a factor \(K\) in each
coordinate. Choose \(\delta<\sigma'/2\), then \(P_0\) sufficiently
large for the fixed mixed orders and this tolerance, and finally
\(T\) large. This proves the lemma, including \(K=1\) and
\(D_i=1\). It prices the \(D\)'s for one specified \(K\); it does
not sum over possible complements.
\end{proof}

\section{Polarization and the cost of a complementary modulus}\label{sec:tariff}

Fix a real parent and an original arithmetic itinerary. We count integers
$u$ in a bin $U/2\le u\le U$, not factorizations chosen before seeing the
parent. Write $G_u=(1/u)\mathbb Z/\mathbb Z$. Suppose that $B_u\subset
G_u$ is nonempty, exactly invariant under multiplication by $3$, and
$B_u+B_u$ is covered by at most $U^\theta$ translates of $B_u$.
The argument also applies at the terminal node, with $u$ the residual
modulus. All entropies and logarithms in this section are in nats.

\begin{lemma}[Polarization in prime digits]\label{lem:digits}
Fix $0<c,\delta,\rho<1$. There exist integers $m,R$, depending only on
these parameters, with the following property. Let $3\nmid u$, let
$\varnothing\ne B\subset G_u$ be invariant under multiplication by $3$,
and suppose $B+B$ is covered by $M$ translates of $B$. Use the Chinese
remainder isomorphism from $G_u\simeq\mathbb Z/u\mathbb Z$ to its
prime-power coordinates. In any fixed order of their base-prime digits
that places lower digits before higher digits in each coordinate, let
$h_i$ be the entropy of digit $i$ of a uniform $X\in B$, conditional on
all preceding digits. Denote its base by $r_i$. For any selection $I$
of digits satisfying
\[
 r_i>R,\qquad \operatorname{ord}_{r_i}(3)\ge r_i^c,
 \qquad h_i\ge\delta\log r_i ,
\]
one has
\[
 \sum_{i\in I}(\log r_i-h_i)
 \le \rho\sum_{i\in I}\log r_i+(m-1)\log M .
\]
\end{lemma}

\begin{proof}
Take $m$ independent copies $X_1,\ldots,X_m$ of $X$, and, independently,
choose $t_1,\ldots,t_m$ uniformly and independently modulo
$\operatorname{ord}_u(3)$, using least nonnegative representatives.
The case $u=1$ is empty. Put $Y=\sum_l3^{t_l}X_l$, and let
$e_i=H(Y_i\mid Y_{<i},\boldsymbol t)$. In bounding $e_i$ below we may
instead condition on $\boldsymbol t$ and the $m$ individual input
histories $X_{l,<i}$: these determine every preceding output digit.
Indeed, for a digit at level $j$ (the units digit is level $0$) of
an $r$-primary coordinate, let $b_l$ be its known residue modulo $r^j$
in $[0,r^j)$, and $v_l$ its current digit. The output digit is
\[
 \left\lfloor r^{-j}\sum_l3^{t_l}b_l\right\rfloor+
 \sum_l(3^{t_l}\bmod r)v_l \pmod r .
\]
The first term is known. The remaining terms are independent conditional
on the individual histories and the times. Adding independent variables
on a finite group cannot lower the entropy of any one summand
(condition on the other summands). Thus $e_i\ge h_i$ for every $i$.

We claim also $e_i\ge(1-\rho)\log r$ whenever the hypotheses for $I$
hold. Write $\mu_z$ for the conditional digit law given a single history
$z$, and $a_z$ for its mass on atoms of size at most $r^{-\delta/4}$.
Separating these atoms from their complement gives
\[
 H(\mu_z)\le\log 2+(\delta/4)\log r+a_z\log r
\]
since the complement has at most $r^{\delta/4}$ elements.
Hence $\mathbb E a_z\ge\delta/2$ for large $r$. Flag a copy if its
current atom is in that small-atom set and $a_z\ge r^{-\delta/8}$.
The probability of a flag is at least $\delta/3$, and conditional on
its history and positive flag a copy has maximum atom at most
$r^{-\delta/8}$. The flags of the copies are independent and do not
depend on the times. We may lower bound the output entropy by
additionally conditioning on all these flags.

We use the prime-field subgroup exponential-sum estimate of
Bourgain, Glibichuk and Konyagin \cite{BGK}: for some $b=b(c)>0$ and
$C=C(c)$, if $\Gamma\subset\mathbb F_r^\times$ is a subgroup of size
at least $r^c$, the absolute value of the average of any nontrivial
additive character over $\Gamma$ is at most $C r^{-b}$.
For a flagged conditional law $\nu$ and $v\ne0$, expand the squared
Fourier coefficient with two independent samples from $\nu$.
The diagonal has mass at most $r^{-\delta/8}$; off the diagonal
$3^{t_l}\bmod r$ is uniform on $\langle3\rangle$, so
\[
 \mathbb E_{t_l}|\widehat\nu(3^{t_l}v)|^2
 \le r^{-\delta/8}+Cr^{-b}\le r^{-b'}
\]
for some $b'>0$ and large $r$. Fix $J$ with $Jb'\ge2$.
Conditional on histories and flags with at least $J$ positive flags,
independence of the times and Parseval give mean conditional collision
probability for the output digit at most
$r^{-1}(1+(r-1)r^{-Jb'})$. The known shift can depend on all the times
without changing this probability; the remaining, possibly unflagged,
Fourier factors have modulus at most one. Entropy is at least minus the
logarithm of collision probability, so Jensen gives mean conditional
entropy at least $\log r-\log(1+r^{-1})$ in this case. Choose $m$ so that
the probability of fewer than $J$ successes in $m$ Bernoulli trials of
parameter $\delta/3$ is at most $\rho/2$, then take $R$ sufficiently
large. This proves the claim. Only independence \emph{between copies}
has been used; the clocks of different primes in one copy may coincide.

Finally $Y$ belongs to the $m$-fold sumset of $B$. Iterating the
translate cover shows that this sumset has at most $M^{m-1}|B|$
elements. Consequently
\[
 0\le\sum_i(e_i-h_i)=H(Y\mid\boldsymbol t)-\log|B|
 \le (m-1)\log M .
\]
Applying the two lower bounds for $e_i$ proves the lemma.
\end{proof}

\begin{lemma}[Prime-base catalogue]\label{lem:basecatalog}
Fix $P_0$, $\xi>0$, and $1/2<s<1$. For $d$ write $d_{\rm big}$ for the
factor composed of prime powers at primes exceeding $P_0$. Consider
either all primes $r\ne3$ for which $d=\operatorname{ord}_r(3)$
satisfies $d_{\rm big}=1$, or, for one fixed $L>1$ having only primes exceeding
$P_0$, all those for which $d_{\rm big}\mid L$ and $d_{\rm big}\ge L^\xi$.
In either catalogue the number of squarefree products of the allowed
bases not exceeding $V\ge1$ is $O_{s,\xi,P_0}(V^s)$, uniformly in $L$.
\end{lemma}
\begin{proof}
At a fixed order $d$ there are $O(d)$ possible primes since their
product divides $3^d-1$. Each is $1\bmod d$. The $j$th smallest is
at least $1+jd$, whence
\[
 \sum_{\operatorname{ord}_r(3)=d}r^{-s}
 \ll_s d^{-s}\sum_{j\ll d}j^{-s}\ll_s d^{1-2s}.
\]
Write $a=2s-1>0$. The sum of $w^{-a}$ over integers with all prime
factors $\le P_0$ is $\prod_{p\le P_0}(1-p^{-a})^{-1}$, with the
product taken over primes. In the second case this constant is
multiplied by at most $\tau(L)L^{-a\xi}=O_{a,\xi}(1)$, by the divisor
bound. Thus $\sum_r r^{-s}$ is uniformly bounded and so is
$\prod_r(1+r^{-s})$. Multiplying the latter by $V^s$ gives the count.
\end{proof}

\begin{proposition}[Adaptive complementary count]\label{prop:tariff}
Fix a large-drop node or a terminal node, its original arithmetic
labels (in particular $n=\operatorname{ord}_q(3)=T^{o(1)}$), and
its actual preceding observations. Take $T^\eta\le U\ll T$ with
fixed $\eta>0$. For every $\sigma>0$, with the covering exponent
$\theta$ sufficiently small and the tolerance in Lemma~\ref{lem:lcm}
sufficiently small, the number of eligible moduli $u\in[U/2,U]$ with
$|B_u|\ge U^\beta$ satisfies, uniformly for $0\le\beta\le1$,
\[
 \#\{u\}\ \ll\ T^{\xi+o(1)}U^{(1-\beta)/2+\sigma}.
\]
At a terminal node the factor $T^\xi$ is omitted. Here $P_0$ may be
chosen after the covering and quotient parameters; the implied
constant may depend on its final value.
\end{proposition}

\begin{proof}
Recall the arithmetic information used at a single node. At a prime
drop the part of the entering lag supported on primes $>P_0$ is $L$.
Outside primes dividing $n$, the low-order part lost has product at
most $T^\xi$ by the itinerary choice; all other lost bases are in
the second catalogue of Lemma~\ref{lem:basecatalog}. At the terminal
all orders divide the terminal lag, so the first catalogue applies.
These are statements about the original itinerary. Indeed, if
$r^e\mid q$ and $r\nmid 3n$, the order at $r^e$ equals the order at
$r$: their quotient is an $r$-power and divides $n$. Such a component
leaves the lag gcd in its entirety or not at all. Partial losses over
$r\mid n$ need not satisfy this assertion.

Choose a large integer $A$, small $c,\delta,\rho>0$ and
$1/2<s<1$ close to $1/2$. Let $m,R$ be as in Lemma~\ref{lem:digits}.
Put aside, in the following order, the parts of $u$ listed here.
Their products, taken separately, have the indicated number of
possible integer names; their heights need not themselves be small.
\begin{enumerate}
\item The part on primes dividing $n$ or at most $R$: $T^{o(1)}$ names.
For each fixed $v>0$ the number is at most
$U^v\prod_{r\mid n\ \mathrm{or}\ r\le R}(1-r^{-v})^{-1}$.
The logarithm of the product is $O_{v,R}(1+\log(2n))=o(\log T)$.
Taking arbitrarily small fixed $v$ gives the assertion.
\item The remaining low-order loss of the itinerary: at most $T^\xi$
names by its height, and no such part is needed at the terminal.
\item All remaining full powers $r^e$ with $e\ge A$: $O_A(U^{2/A})$
names, since
$\prod_r(1+\sum_{e\ge A}r^{-2e/A})<\infty$.
\item All remaining bases with $\operatorname{ord}_r(3)<r^c$,
including their powers: $O_c(U^{3c})$ names. Indeed, at most
$O(Y^{2c})$ such bases lie below $Y$, by summing $O(d)$ over
$d\le Y^c$. Dyadic summation gives $\sum_r r^{-3c}<\infty$
over these bases, so the full geometric Euler product converges.
\end{enumerate}
For the rest, call a full block $r^e$ active if \emph{every} digit
has $h_i\ge\delta\log r$ for $X$ uniform on $B_u$, using a fixed
order for that $u$. Put the product of these whole blocks in $D$, and
put $K=u/D$. Empty products are allowed. The complementary factor
$K$ includes all parts put aside. Lemma~\ref{lem:digits} and
$M\le U^\theta$ imply that the sum of the entropy deficits of the
active digits is at most
$\alpha\log U$, where $\alpha=\rho+(m-1)\theta$.
Conditioning only on previous active digits instead of on the full
histories can only increase the corresponding entropies. The
projection of $B_u$ on its active blocks therefore has at least
$D U^{-\alpha}$ elements. Since $(D,K)=1$, multiplication by $K$
kills exactly the complement and is invertible on the active blocks.
In particular
\[
                 |K\cdot B_u|\ge D U^{-\alpha}.
\]
For each \emph{fixed} integer $K$, use Lemma~\ref{lem:lcm} with
tolerance $\lambda>0$: its sum of $|K\cdot B_{KD}|$ over eligible
$D$ is $\ll (U/K)U^\lambda T^{o(1)}$. Since $D\ge U/(2K)$, there
are at most $O(U^{\alpha+\lambda}T^{o(1)})$ choices of $D$ in
our selection, including when $K=1$ or $D=1$.

It remains to count $K$, which depends on the set $B_u$. Every
remaining inactive block has a digit with deficit at least
$(1-\delta)\log r$. Since all deficits are nonnegative and sum
to $\log u-\log|B_u|$, the product of these distinct bases satisfies
\[
    \sum_{r\ {\rm remaining\ inactive}}\log r
           \le \frac{(1-\beta)\log U}{1-\delta}.
\]
For $1\le e<A$ put $V_e$ equal to the least power of two not below
the product of these bases of exponent $e$ (use $1$ for an empty
product). There are $O_A((1+\log U)^A)$ vectors of bins. In any
nonempty vector,
$\sum_e\log V_e\le A\log2+(1-\beta)\log U/(1-\delta)$.
For each $e$ the eligible base products are squarefree subproducts
of the \emph{original} node catalogue, so at most
$O_{s,\xi,P_0}(V_e^s)$ can occur. Raising a chosen product to the
known exponent $e$ recovers that part of $K$. This counting argument
uses neither a newly valid itinerary after deletion nor independence
among the factors; the bins are a device for a union bound at this
parent, not extra labels on the stationary point distribution.

Multiplication of the counts for the parts put aside, the remaining
inactive names, and the choices at fixed $K$ gives the asserted
estimate with exponent on $U$ at most
\[
 \frac{s(1-\beta)}{1-\delta}
       +\frac2A+3c+\rho+(m-1)\theta+\lambda ,
\]
besides $T^{o(1)}$ and the possible $T^\xi$. First choose
$A,c,\delta,\rho,s$ to make the excess over $(1-\beta)/2$
small uniformly in $\beta$; next fix $m,R$, then $\theta$ and
$\lambda$. The large time-prime cutoff needed to realize the
cover and Lemma~\ref{lem:lcm} can be taken afterwards.
\end{proof}

\section{Assembly of the count}\label{sec:assembly}

\begin{proof}[Proof of Theorem~\ref{thm:main}]
We first count the nonendpoint points with denominator coprime to \(3\)
in one shell \(T/2<q\le T\). Fix the parameters of
Lemma~\ref{lem:itinerary} and one nonempty label class. Choose \(X\)
uniform from it. Let \(E\) and the real traces \(W_i\) be as in
Lemma~\ref{lem:small-traces}. Since the data \((E,W_1,\ldots,W_J,X)\)
include \(X\) and are functions of \(X\), the logarithm of the size of the
label class is
\[
 H(E)+\sum_{i=1}^J H(W_i\mid E,W_1,\ldots,W_{i-1})
       +H(X\mid E,W_1,\ldots,W_J).
\]
The side information and the small steps together cost
\(O(\eta\log(2/\eta))\log T+O_\eta(1)\). We now bound the other
conditional entropies uniformly in the conditioning values that occur.

At a large step, use its dyadic upper bin endpoint \(U\). For each
\(u\) in the bin, the entire possible child trace lies above the set
\(A_u\) of Section~\ref{sec:boolean} with at most two lifts per point.
In this application use all eligible witnesses with the fixed past
and label; any further restrictions can only decrease the counts.
Discard empty fibers and put the nonempty ones into intervals of
width at most \(\Delta>0\) according to
\(\log |B_u|/\log U\). If the lower end of such an interval is \(\beta\),
then Proposition~\ref{prop:tariff} bounds the number of moduli by
\(T^{\xi+o(1)}U^{(1-\beta)/2+\sigma}\). By
\eqref{eq:bool-container}, each contributes at most
\(T^{o(1)}U^{\kappa(\beta+\Delta)+\kappa/(P_0-1)}\) points. Since
\(\kappa>1/2\), their total over all intervals, including the lifts, is
\[
 \ll_{\Delta} T^{\xi+o(1)}
      U^{\kappa+\sigma+\kappa\Delta+\kappa/(P_0-1)} .
\]
At the terminal step, for \(U\ge T^\eta\) use the same argument with
the bound \eqref{eq:bool-terminal}, and omit \(T^\xi\) and the
\(P_0\)-dependent power of \(U\). If the terminal upper bin endpoint
is less than \(T^\eta\), there are at most \(O(U^2)\) possibilities:
there are at most \(U\) moduli and at most \(U\) points per affine
grid. Each conditional entropy is at most the logarithm of the
corresponding support bound.

There are \(O(1/\eta)\) steps. The product of the true large increments
and the terminal modulus is at most \(q\), so the sum of the logarithms
of the corresponding bin upper endpoints is at most
\(\log T+O(1/\eta)\). The total extra loss outside the powers of the bin
endpoints is at most
\(O(\eta\log(2/\eta)+\eta+\xi/\eta)\log T+o(\log T)\), at fixed
parameters. To obtain any prescribed positive slack, first choose
\(\eta\) small and set \(\xi=\eta^3\). Then choose \(\sigma\) and
\(\Delta\) small. Follow the choice of parameters in the proof of
Proposition~\ref{prop:tariff}: fix the digit-entropy tolerances and
number of copies before the covering exponent; next use a fixed
finite covering menu and a fixed moment order for
Lemma~\ref{lem:lcm}. Finally choose \(P_0\) large enough for the
resulting finite number of inputs in Theorem~\ref{thm:native}, as well
as for the loss \(\kappa/(P_0-1)\). Only after these choices let \(T\)
grow. The stationary label count and the rotation losses are then
\(T^{o(1)}\); no claim of stationarity under conditioning is used.
This proves the bound with arbitrary positive slack for one shell.
Geometric summation gives the same bound for all purely periodic
points up to height \(T\), including the two endpoints at \(q=1\).

For a general reduced denominator \(q=3^j d\), with \(3\nmid d\),
deleting the first \(j\) ternary digits gives a point of reduced
denominator \(d\) and leaves at most \(2^j\) possible prefixes.
The integer tails \(0,1\) obey the same upper count. If the pure count
is \(O_\varepsilon(Y^{\kappa+\varepsilon})\) at height \(Y\ge1\),
the full count is bounded by a constant times
\[
 \sum_{0\le j\le\log_3 T} 2^j
                  (T/3^j)^{\kappa+\varepsilon}
 \le T^{\kappa+\varepsilon}\sum_{j\ge0}3^{-j\varepsilon}.
\]
This proves the theorem.
\end{proof}

The preperiod decomposition gives an exact identity. If \(P(Y)\)
counts the points up to reduced height \(Y\) with \(3\nmid q\), then
\begin{equation}\label{eq:preperiod}
 N_C(T)=P(T)+\sum_{j=1}^{\lfloor\log_3 T\rfloor}2^{j-1}P(T/3^j).
\end{equation}
In fact, write a pure tail as \(a/d=0.\overline B_3\), with
\(b\in\{0,2\}\) the last trit of \(B\). This includes the constant
tails \(0,1\). Modulo \(3\) one has \(a/d=-b\), by using the displayed
periodic denominator \(3^{|B|}-1\). Attach a prefix of length \(j\)
and integer value \(A\). The new numerator \(dA+a\) is coprime
to \(d\), and it is prime to \(3\) exactly when the last trit of
\(A\) differs from \(b\). Exactly half the prefixes thus give reduced
denominator \(3^j d\), and every point of this exact \(3\)-adic
valuation is recovered by stripping \(j\) trits.

In particular \(P(Y)\ll Y^\kappa\) would give the matching upper
bound \(N_C(T)\ll T^\kappa\log(2T)\). In fact, at \(T=3^m\) put
\(F_m=2^{-m}P(3^m)\). Identity~\eqref{eq:preperiod} reads
\(2^{-m}N_C(3^m)=F_m+\tfrac12\sum_{k<m}F_k\).
Consequently bounded averages of the nonnegative \(F_m\) are
equivalent to the matching logarithmic upper bound, by monotonicity
between triadic heights. The estimates above do not imply such
bounded averages.

\section{A superlogarithmic lower bound}\label{sec:lower}

We now give a lower bound that does not use
Theorem~\ref{thm:main}. The construction uses many primes that can
cancel simultaneously in the symbolic denominator. The key is to
arrange these cancellations at enough different period lengths, not
just at one height.

\begin{proof}[Proof of Theorem~\ref{thm:lower}]
Write \(F_m=2^{-m}P(3^m)\), where \(P(Y)\) is the pure count defined
in Section~\ref{sec:assembly}. By \eqref{eq:normalized-preperiod}, to bound even \(N_C(T)-P(T)\) from below
it suffices to bound the average of the \(F_m\) from below.

\smallskip
\noindent\emph{A common multiple of many prime orders.}
Fix \(\beta>15/(32\sqrt e)\) with \(\beta<\kappa\).
By Lichtman~\cite[Theorem~1.1]{Lichtman2022}, extending
Baker and Harman~\cite[Theorem~1]{BH1998}, for every sufficiently large
integer \(x\) there are \(L\gg_\beta x/(\log x)^{C_\beta}\) primes
in \((x,2x]\) whose predecessors have no prime factor exceeding
\(x^\beta\), for a fixed \(C_\beta>1\). Take \(\mathcal L\) to be the
set of these primes and put \(Y=x\). Write
\[
 Q=\prod_{p\in\mathcal L}p,\qquad
 R=x^4\prod_{\substack{r\le x^\beta\\r\ {\rm prime}}}
                       r^{\lfloor\log(2x)/\log r\rfloor}.
\]
Every \(p-1\) with \(p\in\mathcal L\) divides \(R\): each of its
prime powers is at most \(2x\), and its prime base is at most
\(x^\beta\). For large \(x\) all the selected primes exceed \(3\).
Fermat's little theorem then gives \(Q\mid 3^n-1\) whenever
\(n\) is a positive multiple of \(R\).
We have \(\log Q\le 2x\log(2x)\) and
\(\log R\ll x^\beta\), the latter from
\(\log R\le4\log x+\pi(x^\beta)\log(2x)\).
Here the elementary Chebyshev bound \(\pi(z)\ll z/\log z\) suffices:
\(\prod_{m<r\le2m, r\ {\rm prime}}r\) divides \(\binom{2m}{m}\),
so dyadic summation gives \(\sum_{r\le z, r\ {\rm prime}}\log r\ll z\).
Separating primes at \(\sqrt z\) gives the claimed bound for \(\pi(z)\).

\smallskip
\noindent\emph{Many cancelled words at each selected period.}
Set
\[
 k=\left\lfloor\frac{Y^\kappa L}{8Y}\right\rfloor,
 \qquad j=\lfloor k\log_3Y\rfloor .
\]
Since \(Y=x\), we have
\(k\gg x^\kappa/(\log x)^{C_\beta}\) for large \(x\). Also \(L\le2Y\),
so \(k\ll x^\kappa\), \(1\le k<L\), and \(j\ll x^\kappa\log x<R\).

For each positive multiple \(n\) of \(R\) let
\[
 X_n=\sum_{t=0}^{n-1}b_t3^t,
\]
where the \(b_t\) are independent fair bits. Thus \(2X_n\) is
uniform among the integers specified by the \(2^n\) length-\(n\)
Cantor words. To compare \(X_n\bmod Q\) with the uniform residue,
consider a nonzero frequency \(s\bmod Q\). Its Fourier coefficient
has modulus
\[
 \left|\mathbb E\exp(2\pi i sX_n/Q)\right|
   =\prod_{t=0}^{n-1}|\cos(\pi s3^t/Q)|.
\]
Let \(\vartheta=\cos(\pi/6)<1\) and
\(B_Q=\lceil\log_3 Q\rceil+2\).
Every block of \(B_Q\) consecutive factors contains one factor at
most \(\vartheta\). Indeed the distance from \(s3^t/Q\) to the
nearest integer is at least \(1/Q\), since \(3\) is invertible
modulo \(Q\) and \(s\ne0\). Whenever that distance is less than
\(1/6\), tripling the phase triples the distance. It cannot stay
below \(1/6\) for the whole block. The Fourier coefficient is
therefore at most \(\vartheta^{\lfloor n/B_Q\rfloor}\).
By character orthogonality, at each single residue the error from
uniform probability is at most this bound. For any residue event
the error is thus at most
\[
 Q\vartheta^{\lfloor n/B_Q\rfloor}
 \le \exp\bigl(-c_0x^3/\log x\bigr)
\]
for an absolute \(c_0>0\) and large \(x\), because \(n\ge R\ge x^4\)
and \(\log Q\le 2x\log(2x)\). This bound is uniform in \(n\).

For a uniform residue modulo \(Q\), the events \(p\mid X_n\),
\(p\in\mathcal L\), are independent by the Chinese remainder
theorem. The probability of exactly \(k\) successes in that model
is at least
\[
 \binom Lk(2Y)^{-k}(1-1/Y)^L
   \gg \left(\frac{L}{2Yk}\right)^k,
\]
since \(\binom Lk\ge(L/k)^k\) and \(L\le2Y\).
This lower bound is at least \(\exp(-O(x^\kappa\log x))\),
using \(L/(2Yk)\ge1/(2Y)\). The uniform Fourier error is
negligible compared with it. It follows that at least
a constant times \(2^n(L/(2Yk))^k\) Cantor words have exactly
\(k\) selected primes dividing \(X_n\).

Each such word represents the point \(2X_n/(3^n-1)\).
The common divisor of its displayed numerator and denominator
is at least \(Y^k\ge3^j\); its actual reduced denominator is
therefore at most \(3^{n-j}\) and is prime to \(3\).
For fixed \(n\) the fractions are distinct because the numerators
at the common denominator \(3^n-1\) are distinct. Shorter least
periods cause no duplication at this fixed \(n\), and we do not
count points from different lengths as disjoint. Thus
\[
 F_{n-j}\gg 2^j\left(\frac{L}{2Yk}\right)^k
      \ge \frac12\left(\frac{Y^\kappa L}{2Yk}\right)^k
      \ge \tfrac12\,4^k .
\]
The indices \(n-j\) form a progression of positive integers with
spacing \(R\), starting at \(R-j\).

\smallskip
\noindent\emph{Averaging over the progression.}
For \(M\ge4R\) there are at least \(M/(2R)\) such indices less
than \(M\). Consequently there is an absolute \(c_1>0\) such that
\[
 \frac1M\sum_{m<M}F_m\gg \frac{4^k}{R}
     \ge\exp\left(c_1\frac{x^\kappa}{(\log x)^{C_\beta}}\right)
\]
for large \(x\), because \(\log R=O(x^\beta)\) and \(\beta<\kappa\).
For each sufficiently large \(M\), take
\(x=\lfloor a(\log M)^{1/\beta}\rfloor\) for a sufficiently small
fixed \(a>0\). The bound \(\log R\ll x^\beta\) ensures \(R\le M^{1/2}\),
so \(M\ge4R\). The lower bound for the average becomes
\[
 \exp\left(c_2\frac{(\log M)^{\kappa/\beta}}
                        {(\log\log M)^{C_\beta}}\right)
\]
up to a constant depending on \(\beta\), with \(c_2>0\). Inserting this in
the preperiod identity bounds \((N_C(3^M)-P(3^M))/(2^M M)\) from below.

For arbitrary large \(T\), put \(M=\lfloor\log_3T\rfloor\).
The numerator \(N_C(T)-P(T)\) counts the points of positive denominator
valuation at \(3\), so it is nondecreasing. Also
\(T^\kappa<2\cdot 2^M\) and \(M\asymp\log T\). Hence the argument gives
constants \(c_\beta>0\) and \(C_\beta>1\) such that
\[
 N_C(T)-P(T)\gg_\beta T^\kappa\log T
 \exp\left(c_\beta\frac{(\log\log T)^{\kappa/\beta}}
                             {(\log\log\log T)^{C_\beta}}\right)
\]
for all sufficiently large \(T\). For every \(0<\rho<\rho_*\), choose
\(\beta>15/(32\sqrt e)\) with \(\beta<\kappa\) and \(\rho<\kappa/\beta\).
The displayed exponent then eventually exceeds \((\log\log T)^\rho\)
and proves the theorem.
\end{proof}

The argument also shows that the averages
\(M^{-1}\sum_{m<M}F_m\) tend to infinity, rather than being large
only along a subsequence. The lower bound does not determine the
optimal subpower factor in \(N_C(T)\).

\begin{proof}[Proof of Theorem~\ref{thm:lower-cylinders}]
Write \(A\) for the integer of the specified prefix. The cylinder map
is \(y\mapsto (A+y)/3^d\), a bijection from \(C\) onto \(I\).
If \(y=a/q\) is reduced and \(3\mid q\), then \(qA+a\) is coprime to
\(q\) and to \(3\), so the image has reduced denominator exactly
\(3^d q\). If instead \(3\nmid q\), the image has denominator
valuation at most \(d\). This proves the identity at every height,
including the endpoints. The lower bound for \(N_C(Y)-P(Y)\) proved
above gives the first inequality. At triadic heights \(T=3^M\)
with \(M\ge d+1\), the same identity reads
\[
 2^{M-d-1}\sum_{0\le m<M-d}F_m .
\]
For \(d\le\theta\lfloor\log_3 T\rfloor\), we have
\(Y\ge T^{1-\theta}\), \(Y^\kappa=2^{-d}T^\kappa\) and
\(\log\log Y=\log\log T+O_\theta(1)\). Apply the first inequality
with any fixed exponent strictly between \(\rho\) and \(\rho_*\)
to absorb the last bounded change of scale. This gives the claimed
uniform estimate.
\end{proof}

\section{Intrinsic approximation}\label{sec:metric}

For \(\tau\ge1\), let \(W_C(\tau)\) consist of the points \(x\in C\)
for which \(|x-a/q|<q^{-\tau}\) for infinitely many distinct reduced
\(a/q\in C\).

\begin{corollary}\label{cor:dimension}
For \(\tau\ge1\), \(\dim_H W_C(\tau)=\kappa/\tau\).
For \(\tau>1\), the subset whose intrinsic approximation exponent
is exactly \(\tau\) has infinite \(\kappa/\tau\)-dimensional Hausdorff
measure in every positive-radius ball in \(C\). Natural Cantor-almost
every point has intrinsic approximation exponent \(1\).
\end{corollary}

\begin{proof}
The exponent of \(x\) is the supremum of the \(s\) for which
\(x\in W_C(s)\). At dimension \(b>\kappa/\tau\), choose
\(\varepsilon<\tau b-\kappa\). Covering by the balls from each
denominator shell \(2^j\le q<2^{j+1}\) costs
\(O(2^{j(\kappa+\varepsilon-\tau b)})\) by
Theorem~\ref{thm:main}. The tail costs tend to zero. In particular,
\(W_C(\tau)\) has natural Cantor measure zero for \(\tau>1\),
since that measure is comparable to \(\mathcal H^\kappa\) on \(C\).

For the lower bound at \(\tau>1\), use balls in \(C\) of radius
\(\frac12\,3^{-n\tau}\) about the left endpoints of the level-\(n\)
Cantor intervals. These endpoints have denominators at most \(3^n\).
Set \(g(r)=r^\kappa\) and \(f(r)=r^{\kappa/\tau}\). Enlarging each
radius \(R\) to \(g^{-1}(f(R))\) gives radius
\(2^{-1/\tau}3^{-n}>3^{-n-1}\). The enlarged limsup has full Cantor
measure: almost every digit sequence visits first children infinitely
often. Since \(C\) is \(\kappa\)-Ahlfors regular, the general mass
transference principle~\cite[Theorem~3]{BV2006} gives full local
\(\mathcal H^f\) measure for the original limsup. This measure is
infinite on each positive-radius ball in \(C\). Membership gives
infinitely many distinct centers except possibly at countably many
centers themselves, which have zero \(\mathcal H^f\) measure.
Removing \(\bigcup_{m\ge1}W_C(\tau+1/m)\) costs zero
\(\mathcal H^f\) by the upper dimension bound, and leaves exactly
exponent \(\tau\). For \(\tau=1\), truncation gives infinitely many
distinct approximants for every irrational point of \(C\); endpoints
of each cylinder have distance at most \(3^{-n}\) and denominator
at most \(3^n\). The preceding null-set result and countable union
exclude every larger exponent almost surely.
\end{proof}

\section*{Acknowledgments}
OpenAI's GPT-6.2 Astra carried out the mathematical discovery and proof
development reported here and wrote the manuscript. The human author
directed the research and its presentation.
\begingroup
\small
\bibliographystyle{abbrv}
\bibliography{references}
\endgroup
\end{document}
